% ============================================================
% Plantilla: Currículum vitae (una página)
% ============================================================
\documentclass[10pt,a4paper]{article}

\usepackage[utf8]{inputenc}
\usepackage[T1]{fontenc}
\usepackage{lmodern}
\usepackage[spanish]{babel}
\usepackage[margin=1.8cm]{geometry}
\usepackage{xcolor}
\usepackage{tabularx}
\usepackage[hidelinks]{hyperref}

\definecolor{acento}{RGB}{19,128,79}
\pagestyle{empty}
\setlength{\parindent}{0pt}

% Título de sección con línea de color
\newcommand{\seccion}[1]{%
  \vspace{10pt}{\large\bfseries\color{acento}#1}\par
  \vspace{2pt}{\color{acento}\rule{\linewidth}{0.8pt}}\par\vspace{4pt}}

% Entrada: \entrada{Puesto o título}{Fechas}{Lugar}{Descripción}
\newcommand{\entrada}[4]{%
  \textbf{#1} \hfill {\small #2}\par
  {\itshape #3}\par
  {\small #4}\par\vspace{6pt}}

\begin{document}

\begin{center}
  {\Huge\bfseries Nombre Apellido}\par\vspace{4pt}
  {\large Ingeniero de Software}\par\vspace{6pt}
  {\small Ciudad, País \;|\; +593 000 000 000 \;|\; \href{mailto:correo@ejemplo.com}{correo@ejemplo.com}
  \;|\; \href{https://github.com/usuario}{github.com/usuario}}
\end{center}

\seccion{Perfil}
Profesional orientado a resultados, con experiencia en desarrollo web, seguridad de la
información y trabajo en equipos ágiles.

\seccion{Experiencia}
\entrada{Desarrollador Full Stack}{2024 -- Actualidad}{Empresa Ejemplo S.A.}{%
  Desarrollo de aplicaciones web con React y Node.js; integración continua y pruebas automatizadas.}
\entrada{Practicante de Desarrollo}{2023}{Institución Ejemplo}{%
  Mantenimiento de sistemas internos y documentación técnica.}

\seccion{Educación}
\entrada{Ingeniería en Software}{2020 -- 2025}{Nombre de la Universidad}{%
  Trabajo de titulación: título del trabajo.}

\seccion{Habilidades}
\begin{tabularx}{\linewidth}{@{}lX@{}}
  \textbf{Lenguajes:} & JavaScript, TypeScript, Python, SQL \\
  \textbf{Herramientas:} & Git, Docker, Supabase, \LaTeX \\
  \textbf{Metodologías:} & Scrum, principios SOLID, pruebas unitarias \\
\end{tabularx}

\seccion{Idiomas}
Español (nativo) \;\textbullet\; Inglés (intermedio B2)

\end{document}
